\documentclass[11pt,a4paper]{article}

% ============================================
% ENCODING & FONTS
% ============================================
\usepackage[utf8]{inputenc}
\usepackage[T1]{fontenc}
\usepackage{lmodern}

% ============================================
% PAGE LAYOUT (Golden Ratio)
% ============================================
\usepackage[margin=1.618cm, top=2.618cm, bottom=2.618cm]{geometry}

% ============================================
% ESSENTIAL PACKAGES
% ============================================
\usepackage{float}
\usepackage{caption}
\usepackage{setspace}
\usepackage{fancyhdr}
\usepackage{xcolor}
\usepackage{hyperref}
\usepackage{amsmath}
\usepackage{amssymb}
\usepackage{amsthm}
\usepackage{booktabs}
\usepackage{longtable}
\usepackage{array}
\usepackage{listings}
\usepackage{graphicx}
% Unicode degree symbol (native in XeTeX/LuaTeX)
\ifdefined\DeclareUnicodeCharacter
  \DeclareUnicodeCharacter{00B0}{\ensuremath{^\circ}}
\fi

% ============================================
% LISTINGS CONFIGURATION
% ============================================
\lstset{
    basicstyle=\small\ttfamily,
    breaklines=true,
    frame=single,
    keepspaces=true,
    showstringspaces=false,
    breakatwhitespace=true,
    aboveskip=0.8em,
    belowskip=0.8em
}

% ============================================
% TITLE FORMATTING
% ============================================
\usepackage{titling}
\pretitle{\LARGE\bfseries}
\posttitle{\vspace{-0.4em}}
\preauthor{}
\postauthor{}
\predate{}
\postdate{}
\setlength{\droptitle}{-2.0em}

% ============================================
% HEADER/FOOTER
% ============================================
\setlength{\headheight}{14pt}
\pagestyle{fancy}
\fancyhf{}
\fancyhead[L]{Explicit $\Gtwo$ Metric on a Compact TCS 7-Manifold}
\fancyhead[R]{\thepage}
\renewcommand{\headrulewidth}{0.2pt}

% ============================================
% HYPERREF
% ============================================
\hypersetup{
    colorlinks=true,
    linkcolor=blue,
    citecolor=blue,
    urlcolor=blue,
    pdftitle={An Explicit Approximate G2 Metric on a Compact TCS 7-Manifold with Certified Torsion-Free Completion},
    pdfauthor={Brieuc de La Fourniere}
}

% ============================================
% SPACING
% ============================================
\setstretch{1.2}
\setlength{\parskip}{0.4em}
\setlength{\parindent}{0pt}

% ============================================
% CUSTOM COMMANDS
% ============================================
\newcommand{\E}{\mathrm{E}}
\newcommand{\Gtwo}{\mathrm{G}_2}
\newcommand{\Kseven}{K_7}
\newcommand{\dimE}{\mathrm{dim}}
\newcommand{\Weyl}{\mathrm{Weyl}}
\newcommand{\rk}{\mathrm{rank}}
\newcommand{\proven}{\textsc{Proven}}
\newcommand{\topomark}{\textsc{Topological}}
\newcommand{\derived}{\textsc{Derived}}
\newcommand{\AdS}{\mathrm{AdS}}
\newcommand{\SU}{\mathrm{SU}}
\newcommand{\SO}{\mathrm{SO}}
\newcommand{\U}{\mathrm{U}}
\newcommand{\Spin}{\mathrm{Spin}}
\newcommand{\Hol}{\mathrm{Hol}}
\newcommand{\Ric}{\mathrm{Ric}}
\newcommand{\RMS}{\mathrm{RMS}}
\newcommand{\ACyl}{\mathrm{ACyl}}
\newcommand{\PINN}{\mathrm{PINN}}
\newcommand{\TCS}{\mathrm{TCS}}
\newcommand{\Sym}{\mathrm{Sym}}
\newcommand{\chol}{\mathrm{chol}}
\newcommand{\Tr}{\mathrm{Tr}}
\newcommand{\Lip}{\mathrm{Lip}}
\newcommand{\dist}{\mathrm{dist}}
\newcommand{\softplus}{\mathrm{softplus}}

\newcommand{\norm}[1]{\|#1\|}
\newcommand{\abs}[1]{|#1|}

\newtheorem{theorem}{Theorem}[section]
\newtheorem{proposition}[theorem]{Proposition}
\theoremstyle{remark}
\newtheorem*{remark}{Remark}

\pdfstringdefDisableCommands{%
  \def\Gtwo{G2}%
  \def\Kseven{K7}%
  \def\E{E}%
  \def\dimE{dim}%
  \def\Weyl{Weyl}%
  \def\rk{rank}%
  \def\proven{Proven}%
  \def\topomark{Topological}%
  \def\AdS{AdS}%
  \def\SO{SO}%
  \def\Spin{Spin}%
}

% ============================================
% DOCUMENT
% ============================================

\title{%
\LARGE\textbf{An Explicit Approximate $\Gtwo$ Metric on a Compact\\[0.2em]TCS 7-Manifold with Certified Torsion-Free Completion}
}
\author{}
\date{}

\begin{document}

\maketitle
\noindent\rule{\textwidth}{0.2pt}

\noindent\textbf{Author}: Brieuc de La Fourni\`ere

\noindent Independent researcher

\vfill

\begin{abstract}
We construct an explicit approximate Riemannian metric on a compact 7-manifold of
twisted connected sum (TCS) type, with Betti numbers $b_2 = 21$ and $b_3 = 77$. The metric
is given in closed form as a Chebyshev--Cholesky expansion with 169 parameters, defined
on a one-dimensional seam ansatz (reduction along the neck coordinate~$s$). All
certification is performed within this reduced ansatz; the full 7D fiber-dependent
torsion-free equation is not solved. A computational Newton--Kantorovich (NK) certificate
establishes the existence of a unique torsion-free $\Gtwo$ metric $g^*$ within a relative
distance of $4.86 \times 10^{-6}$ of our explicit iterate, which itself has residual torsion
$\norm{T}_{C^0} = 2.98 \times 10^{-5}$ (a factor ${\sim}3000$ reduction from the initial approximation
in 5 Joyce-type iterations). All NK constants (including the inverse bound~$\beta$,
the residual~$\eta$, and the Lipschitz bound~$\omega$ on the Fr\'echet derivative) are computed
with conservative safety factors, yielding the NK contraction parameter
$h = 6.65 \times 10^{-8}$, a margin of ${\times}7.5$~million below the threshold. Full interval
verification of operator norms is left for future work. The K3 fiber variation is
bounded over all 220,000 sample points and contributes at most 0.07\% to the torsion.
The certification chain further includes structural verification (SPD by Cholesky
construction, $\det = 65/32$ algebraically) and a holonomy proof ($\Hol(g^*) = \Gtwo$ via
Joyce~\cite{joyce2000}, Proposition~11.2.3). Separate bounds on $\norm{d\varphi}$ and $\norm{d{*}\varphi}$ confirm that
the small torsion is not due to cancellation. The metric, reconstruction algorithm,
and a self-contained companion notebook ($< 1$~minute runtime) for independent
verification are provided as supplementary material.
\end{abstract}

\vfill
\noindent\rule{\textwidth}{0.2pt}

\newpage
\tableofcontents

\newpage

% ============================================
\section{Introduction}\label{sec:intro}
% ============================================

\subsection{Context}\label{sec:context}

A compact Riemannian 7-manifold $(M^7, g)$ has holonomy contained in the exceptional
Lie group $\Gtwo \subset \SO(7)$ if and only if it admits a torsion-free $\Gtwo$-structure, i.e., a
closed and coclosed 3-form $\varphi \in \Omega^3(M)$~\cite{harvey1982, bryant1987}. Full holonomy $\Gtwo$ (as opposed to a
proper subgroup) requires additionally that $M$ be simply connected and not a
Riemannian product.

Joyce~\cite{joyce1996, joyce2000} proved the existence of compact examples by resolving singularities of
$T^7/\Gamma$ orbifolds. Kovalev~\cite{kovalev2003} introduced the twisted connected sum (TCS) construction,
and Corti--Haskins--Nordstr\"om--Pacini~\cite{chnp2015} systematized the method and classified many
topological types. These results establish the existence of $\Gtwo$ metrics to within a
small (controlled) error of an approximate solution, but do not yield pointwise
numerical values of the metric tensor. Substantial numerical work exists for
\textit{non-compact} $\Gtwo$ manifolds (see e.g.~\cite{brandhuber2001}), but to our knowledge, no explicit metric
$g_{ij}(x)$ has been tabulated for a compact example.

\subsection{Objective}\label{sec:objective}

This paper presents an explicit metric on one particular compact TCS manifold,
together with a certification chain guaranteeing the existence of a nearby
torsion-free metric with holonomy exactly $\Gtwo$. The construction proceeds in three
stages:

\begin{enumerate}
    \item A Chebyshev--Cholesky parametrization of a one-parameter family of $7 \times 7$ metrics
          $g(s)$ along the TCS neck coordinate (the \textit{seam ansatz}).
    \item A Gauss--Newton iteration that reduces the torsion by a factor of ${\sim}3000$ in 5 steps.
    \item A computational Newton--Kantorovich certificate that bounds the distance to the
          exact torsion-free solution.
\end{enumerate}

\subsection{Scope and claims}\label{sec:scope}

To avoid ambiguity, we distinguish four levels of assertion in this paper:

\begin{table}[H]
\centering
\small
\begin{tabular}{p{3.2cm}p{5.5cm}p{7cm}}
\toprule
\textbf{Level} & \textbf{What is claimed} & \textbf{How it is established} \\
\midrule
\textbf{Algebraic} & $\det(g) = 65/32$ exactly; $g$ SPD everywhere; $|\varphi|^2 = 42$ & By construction: Cholesky + softplus + det normalization (\S\ref{sec:cheb_chol}) \\[0.5em]
\textbf{Comp.\ certified} & A unique torsion-free $\Gtwo$ metric $g^*$ exists within $\dist \leq 4.86 \times 10^{-6}$ of $g_5$ & Computational NK certificate with conservative safety factors (\S\ref{sec:nk_cert}) \\[0.5em]
\textbf{Comp.\ certified} & $\Hol(g^*) = \Gtwo$ & Joyce~\cite{joyce2000}, Prop.~11.2.3 + topological data (\S\ref{sec:holonomy}) \\[0.5em]
\textbf{Comp.\ certified} & K3 fiber impact $\leq 0.07\%$ of 1D torsion & Fr\'echet derivative + sup over 220k K3 points (\S\ref{sec:k3_fiber}) \\[0.5em]
\textbf{Computed} & The explicit iterate $g_5$ has $\norm{T}_{C^0} = 2.98 \times 10^{-5}$ & Spectral differentiation on Chebyshev grid (\S\ref{sec:torsion}) \\[0.5em]
\textbf{Computed} & All geometric invariants to 14 digits & IEEE 754 float64 arithmetic (\S\ref{sec:invariants}) \\[0.5em]
\textbf{Not claimed} & Full 7D fiber-dependent torsion-free solution & All certification within the 1D seam ansatz $g(s)$ (\S\ref{sec:model_hier}, \S\ref{sec:limitations}) \\[0.5em]
\textbf{Not claimed} & Extension to the ACyl bulk beyond the exponential tail model & See \S\ref{sec:acyl} \\
\bottomrule
\end{tabular}
\end{table}

We emphasize the distinction between \textit{algebraic} properties (which hold rigorously by
construction) and \textit{computationally certified} properties (which rely on spectral
discretization for $\beta$, finite-difference Jacobians for $\omega$, and IEEE~754 arithmetic
throughout). Full interval verification of operator norms, in the sense of
computer-assisted proofs, is left for future work (\S\ref{sec:future}).

\subsection{Outline}\label{sec:outline}

Section~\ref{sec:manifold} describes the TCS manifold and its topology. Section~\ref{sec:metric} presents the metric
parametrization and model hierarchy. Section~\ref{sec:norms} defines the norms used throughout.
Section~\ref{sec:torsion} contains the torsion analysis. Section~\ref{sec:certification} gives the certification results.
Section~\ref{sec:invariants} summarizes the geometric invariants. Section~\ref{sec:discussion} discusses limitations,
related work, and future directions. Section~\ref{sec:reproducibility} describes the supplementary materials.
Appendices provide the baseline metric, K3 fiber data, Chebyshev coefficients, and a
remark connecting the input data to a related physics framework.

% ============================================
\section{The Manifold}\label{sec:manifold}
% ============================================

\subsection{Twisted Connected Sum Construction}\label{sec:tcs}

The compact 7-manifold $\Kseven$ is built following the Kovalev TCS construction~\cite{kovalev2003, chnp2015}. Two
asymptotically cylindrical (ACyl) Calabi--Yau threefolds $M_1$, $M_2$ are crossed with $S^1$
and glued along a common neck region diffeomorphic to $\mathrm{K3} \times T^2 \times I$, where $I = [0, 1]$
is the seam interval:
\begin{equation}\label{eq:tcs_gluing}
\Kseven = (M_1 \times S^1) \;\cup_{\mathrm{K3} \times T^2 \times I}\; (M_2 \times S^1)
\end{equation}

\begin{figure}[ht]
\centering
\includegraphics[width=0.65\textwidth]{figures/K7_TCS_light.pdf}
\caption{Three-dimensional visualization of $\Kseven$ as a surface of revolution, with
color encoding the torsion intensity $\mathcal{T}(s)$. The two bulbous caps represent the ACyl
bulks $M_-$ and $M_+$ (low torsion, teal); the narrow neck carries 97\% of the torsion
(orange-red). The radial profile reflects the metric eigenvalue structure. Key
invariants annotated at bottom.}
\label{fig:K7_3d}
\end{figure}

\textbf{Building blocks.} The construction uses a pair of semi-Fano threefolds $(Z_1, S_1)$
and $(Z_2, S_2)$, where $S_i \in |{-}K_{Z_i}|$ are smooth K3 divisors. The ACyl Calabi--Yau
threefolds $M_i = Z_i \setminus S_i$ serve as the two halves:

\begin{table}[H]
\centering
\begin{tabular}{lll}
\toprule
\textbf{Component} & \textbf{Description} & \textbf{Topological data} \\
\midrule
$M_1$ & ACyl CY threefold (semi-Fano building block) & $b_2(M_1) = 11$, $b_3(M_1) = 40$ \\
$M_2$ & ACyl CY threefold (semi-Fano building block) & $b_2(M_2) = 10$, $b_3(M_2) = 37$ \\
K3 fiber & Common anticanonical K3 surface & $\chi = 24$, $h^{1,1} = 20$ \\
\bottomrule
\end{tabular}
\end{table}

\subsection{Topology}\label{sec:topology}

We consider a TCS 7-manifold of topological type $(b_2 = 21, b_3 = 77)$. These Betti
numbers are computed from the Mayer--Vietoris sequence for the TCS decomposition,
which involves the matching data of the two building blocks: specifically, the
images of the restriction maps $H^2(M_i) \to H^2(\mathrm{K3})$ and their mutual position in the
K3 lattice $\Lambda_{\mathrm{K3}}$ of signature $(3, 19)$. For the particular building block pair
above, the polarization lattices $N_1 \subset \Lambda_{\mathrm{K3}}$ (rank 11) and $N_2 \subset \Lambda_{\mathrm{K3}}$ (rank 10)
satisfy $N_1 \cap N_2 = \{0\}$, so that the Mayer--Vietoris sequence yields:
\begin{align}
b_2(\Kseven) &= \rk(N_1) + \rk(N_2) = 11 + 10 = 21 \label{eq:b2}\\
b_3(\Kseven) &= b_3(M_1) + b_3(M_2) = 40 + 37 = 77 \label{eq:b3}
\end{align}

The additivity of $b_3$ holds because $H^3(\mathrm{K3}) = 0$; the additivity of $b_2$ holds because
the matching is \textit{orthogonal} ($N_1 \cap N_2 = \{0\}$). For non-orthogonal matchings, correction
terms involving the cokernel of the restriction maps would appear; see~\cite{chnp2015}, \S4 for the
general formula.

Full Betti spectrum: $(1, 0, 21, 77, 77, 21, 0, 1)$. Euler characteristic $\chi(\Kseven) = 0$.

\begin{figure}[ht]
\centering
\includegraphics[width=0.75\textwidth]{figures/fig1_tcs_schematic.pdf}
\caption{Schematic of the TCS decomposition. The three atlas charts $U_L$, $U_N$,
$U_R$ cover the left ACyl bulk ($s \in [-2, 0]$), neck ($s \in [0, 1]$), and right ACyl bulk
($s \in [1, 3]$) respectively. Betti contributions from each building block and the
resulting global topology are indicated.}
\label{fig:tcs_schematic}
\end{figure}

\textbf{Simple connectivity.} Each piece $M_i \times S^1$ has $\pi_1 = \mathbb{Z}$ (from the $S^1$ factor, since
the ACyl Calabi--Yau $M_i$ is simply connected). The TCS gluing map $\Phi$ involves a
hyper-K\"ahler rotation on the K3 fiber that exchanges the two circle factors in the
neck region $\mathrm{K3} \times S^1 \times S^1$: the $S^1$ fiber of one piece is identified with the $S^1$ factor
of the other. By the Seifert--van Kampen theorem applied to the decomposition
$\Kseven = (M_1 \times S^1) \cup (M_2 \times S^1)$ with overlap $\mathrm{K3} \times T^2$, the two $\mathbb{Z}$ generators are
identified across the gluing, yielding $\pi_1(\Kseven) = \{1\}$. In particular $b_1(\Kseven) = 0$.

The moduli count is $b_3 = 77$, consistent with the Joyce--Karigiannis deformation theory
for compact $\Gtwo$ manifolds~\cite{joyce1996, joyce2000}.

% ============================================
\section{The Metric}\label{sec:metric}
% ============================================

\subsection{Model Hierarchy}\label{sec:model_hier}

The full metric on $\Kseven$ would be a field $g_{ij}(s, \theta, y_1, \ldots, y_4, \psi)$ depending on all
7 coordinates. In this paper, we work with a \textbf{reduced seam ansatz}: a one-parameter
family of constant-fiber metrics $g_{ij}(s)$, where $s$ is the neck coordinate and the
metric on each fiber $\mathrm{K3} \times T^2$ at fixed $s$ is taken to be $s$-dependent but
spatially homogeneous on the fiber.

This dimensional reduction is motivated by the TCS geometry: the neck region is
diffeomorphic to $\mathrm{K3} \times T^2 \times [0, 1]$, and the metric is expected to vary primarily
along the gluing direction~$s$.

The relationship between the seam ansatz and the full problem is:

\begin{table}[H]
\centering
\small
\begin{tabular}{p{3.5cm}p{4.5cm}p{7.5cm}}
\toprule
\textbf{Level} & \textbf{Object} & \textbf{Status} \\
\midrule
\textbf{Seam ansatz} & $g_{ij}(s)$, 169 parameters & Fully specified (this paper) \\[0.3em]
\textbf{K3 fiber metric} & $g_{\mathrm{CY}}(y)$ on $\mathrm{CI}(1,2,2,2) \subset \mathbb{P}^6$ & Computed independently via cymyc~\cite{cymyc} (Appendix~\ref{app:k3}) \\[0.3em]
\textbf{Fiber-coupled metric} & $g(s, y) = g(s) + \text{corrections from } g_{\mathrm{CY}}(y)$ & K3 impact certified over 220k points: 0.07\% of torsion bound (\S\ref{sec:k3_fiber}) \\[0.3em]
\textbf{Full 7D metric} & $g(s, \theta, y, \psi)$ & Not constructed; noted as future work \\
\bottomrule
\end{tabular}
\end{table}

The torsion analysis (\S\ref{sec:torsion}) and the NK certification (\S\ref{sec:nk_cert}) are carried out entirely
within the seam ansatz. The K3 verification (\S\ref{sec:k3_fiber}) certifies \textit{a posteriori} that
the K3 fiber variation contributes at most 0.07\% to the torsion bound. However,
the full 7D torsion-free equation on $\Kseven$ has not been solved directly; the
certification establishes convergence within the 1D family $g(s)$.

\subsection{Coordinate System}\label{sec:coordinates}

We use an atlas with three charts covering the TCS manifold:
\begin{itemize}
    \item $U_L$ (left ACyl bulk): $s \in [-2, 0]$
    \item $U_N$ (neck): $s \in [0, 1]$
    \item $U_R$ (right ACyl bulk): $s \in [1, 3]$
\end{itemize}

The coordinate index convention on each fiber at fixed $s$:

\begin{table}[H]
\centering
\begin{tabular}{cll}
\toprule
\textbf{Index} & \textbf{Direction} & \textbf{Type} \\
\midrule
0 & $s$ (seam) & Neck parameter \\
1 & $\theta$ (circle fiber) & $T^2$ \\
2 & $y_1$ (K3 real part 1) & K3 \\
3 & $y_2$ (K3 imaginary part 1) & K3 \\
4 & $y_3$ (K3 real part 2) & K3 \\
5 & $y_4$ (K3 imaginary part 2) & K3 \\
6 & $\psi$ (circle fiber) & $T^2$ \\
\bottomrule
\end{tabular}
\end{table}

K3 directions: $\{2, 3, 4, 5\}$. Torus directions: $\{1, 6\}$.

\subsection{Parametrization: Chebyshev--Cholesky Expansion}\label{sec:cheb_chol}

The metric on the neck ($s \in [0, 1]$) is:
\begin{equation}\label{eq:metric_chol}
g_{ij}(s) = [L(s)\, L(s)^\top]_{ij}
\end{equation}
where $L(s)$ is a $7 \times 7$ lower-triangular matrix expanded in Chebyshev polynomials:
\begin{equation}\label{eq:cheb_expansion}
L_{\mathrm{flat},j}(s) = \sum_{k=0}^{K} c_{kj}\, T_k(2s - 1), \qquad K = 5,\quad j = 0, \ldots, 27
\end{equation}

Here $T_k$ denotes the $k$-th Chebyshev polynomial of the first kind, and the 28 entries
$L_{\mathrm{flat},j}$ are the independent components of the lower triangle of $L$, enumerated in
row-major order:
\[
j\!:\quad \underbrace{L_{00}}_{0}\;\; \underbrace{L_{10}\;L_{11}}_{1,2}\;\; \underbrace{L_{20}\;L_{21}\;L_{22}}_{3,4,5}\;\; \cdots\;\; \underbrace{L_{64}\;L_{65}\;L_{66}}_{25,26,27}
\]

\textbf{Softplus activation on the diagonal.} To guarantee strict positive-definiteness
($g = LL^\top$ with $L$ nonsingular), the diagonal entries of $L$ are passed through a softplus
function:
\begin{equation}\label{eq:softplus}
L_{ii} = \softplus(L_{\mathrm{flat},j}) = \log(1 + \exp(L_{\mathrm{flat},j})) \quad \text{for } j \in \{0, 2, 5, 9, 14, 20, 27\}
\end{equation}

All off-diagonal entries are used directly.

\textbf{Determinant normalization.} At each $s$, the metric is rescaled to enforce a fixed
determinant (see Appendix~\ref{app:related} for the origin of this value):
\begin{equation}\label{eq:det_norm}
\det(g(s)) = 65/32 \qquad \text{(exactly, by construction)}
\end{equation}
via the scaling $L(s) \to \alpha(s) \cdot L(s)$ where $\alpha^7 \det(L) = \sqrt{65/32}$.

\textbf{Total parameter count:} 6 Chebyshev modes $\times$ 28 Cholesky entries = 168 parameters,
plus 1 ACyl decay rate $\gamma$ = \textbf{169 parameters total}.

\begin{figure}[ht]
\centering
\includegraphics[width=0.8\textwidth]{figures/fig3_metric_profile.pdf}
\caption{Eigenvalue profile of $g(s)$ along the full domain $s \in [-2, 3]$. Three
distinct scales are visible: the seam direction $\lambda_0 \approx 6.5$ (dominant), the torus
directions $\lambda_{1,6} \approx 2.9$, and the K3 directions $\lambda_{2\text{--}5} \approx 1.1$. The metric is
nearly constant in the ACyl bulks, with variation concentrated in the neck.
$\det(g) = 65/32$ exactly by construction; condition number $\kappa \leq 3.88$.}
\label{fig:metric_profile}
\end{figure}

\subsection{ACyl Extension}\label{sec:acyl}

For $s$ outside the neck $[0, 1]$, the metric decays exponentially toward the asymptotic
Calabi--Yau cross-section metric:
\begin{equation}\label{eq:acyl_decay}
g(s) \to g_\infty + [g(s_{\mathrm{bdy}}) - g_\infty] \cdot \exp(-2\gamma|s - s_{\mathrm{bdy}}|)
\end{equation}
with decay rate $\gamma = 5.811297$ and bulk domain extending to $s \in [-2, 3]$.
The neck carries 97.2\% of the total torsion; the tails contribute only 2.8\%.

This exponential decay model is an approximation: the true ACyl metric on $M_i$
approaches a Calabi--Yau cylinder metric at a rate controlled by the first eigenvalue
of the Laplacian on the K3 cross-section. The decay rate $\gamma$ is fitted to match the
numerical boundary data and is consistent in order of magnitude with known ACyl
decay estimates~\cite{kovalev2003, chnp2015}.

\subsection{Chebyshev Coefficients}\label{sec:cheb_coeffs}

The $6 \times 28$ coefficient matrix $C = (c_{kj})$ is given in full in the supplementary
data file (see \S\ref{sec:reproducibility} and companion notebook). The spectral structure is:

\begin{table}[H]
\centering
\begin{tabular}{ccl}
\toprule
\textbf{Mode} $k$ & $\max_j |c_{kj}|$ & \textbf{Interpretation} \\
\midrule
0 & 2.559 & Constant term (carries $> 99.99\%$ of metric) \\
1 & $4.0 \times 10^{-3}$ & Linear variation along neck \\
2 & $8.5 \times 10^{-8}$ & Numerical noise (negligible) \\
3 & $5.6 \times 10^{-5}$ & Small cubic correction \\
4 & $1.4 \times 10^{-7}$ & Numerical noise (negligible) \\
5 & $5.2 \times 10^{-5}$ & Small quintic correction \\
\bottomrule
\end{tabular}
\end{table}

\textbf{Effective degrees of freedom:} ${\sim}56$ significant parameters (rows 0--1), plus ${\sim}56$ small
corrections (rows 3, 5). Rows 2, 4 are noise at the $10^{-8}$ level.

\subsection{Reconstruction Algorithm}\label{sec:reconstruction}

Given the coefficient matrix $C$ and decay rate $\gamma$, the metric $g(s)$ at any point $s$ is
reconstructed in 6 steps:

\begin{enumerate}
    \item \textbf{Chebyshev evaluation:} For each $j \in \{0, \ldots, 27\}$, compute $L_{\mathrm{flat},j}(s) = \sum_k c_{kj}\, T_k(2s{-}1)$.
    \item \textbf{Reshape:} Pack the 28 values into a $7 \times 7$ lower-triangular matrix $L(s)$.
    \item \textbf{Softplus diagonal:} Replace $L_{ii} \leftarrow \log(1 + \exp(L_{ii}))$ for the 7 diagonal entries.
    \item \textbf{Metric:} Compute $g(s) = L(s) \cdot L(s)^\top$.
    \item \textbf{Det-normalize:} Rescale $g(s) \leftarrow g(s) \cdot (65/32 \;/\; \det(g(s)))^{1/7}$.
    \item \textbf{ACyl tails:} For $s \notin [0, 1]$, apply exponential decay with rate $\gamma = 5.811297$.
\end{enumerate}

The $\Gtwo$ 3-form $\varphi$ and 4-form $\psi = {*}\varphi$ are then:
\begin{align}
\varphi_{ijk}(s) &= \sum_{(abc) \in \text{Fano}} \mathrm{sign}(abc) \cdot L_{ia}(s)\, L_{jb}(s)\, L_{kc}(s) \label{eq:phi_fano}\\
\psi_{ijkl}(s) &= \sum_{(abcd) \in \text{CoFano}} \mathrm{sign}(abcd) \cdot L_{ia}(s)\, L_{jb}(s)\, L_{kc}(s)\, L_{ld}(s) \label{eq:psi_fano}
\end{align}
where the sums run over the 7 associative triples and 7 coassociative quadruples of
the Fano plane, with standard orientations.

% ============================================
\section{Norm Definitions and Domain}\label{sec:norms}
% ============================================

All norms are defined on the \textbf{full extended neck domain} $s \in [-2, 3]$, evaluated on
the 1D seam ansatz $g(s)$ at Chebyshev collocation nodes ($N = 100$ on the neck, with
additional tail nodes).

\subsection{Metric Distance}\label{sec:metric_dist}

The relative metric perturbation between two metrics $g_0$, $g_1$ is:
\begin{equation}\label{eq:metric_dist}
\delta g / g := \sup_{s \in [-2,3]} \frac{\norm{g_1(s) - g_0(s)}_F}{\norm{g_0(s)}_F}
\end{equation}
where $\norm{\cdot}_F$ denotes the Frobenius norm of the $7 \times 7$ matrix.

\subsection{Torsion Norms}\label{sec:torsion_norms}

The torsion of a $\Gtwo$ structure $(\varphi, \psi)$ is $T = (d\varphi, d{*}\varphi)$. In the 1D seam ansatz, the
only nonvanishing derivatives are $\partial_s \varphi$ and $\partial_s \psi$. The proper (metric-contracted)
torsion components are:
\begin{align}
|d\varphi|^2(s) &= 4 \cdot g^{00}(s) \cdot g^{aa'}(s)\, g^{bb'}(s)\, g^{cc'}(s) \cdot (\partial_s \varphi_{abc})(\partial_s \varphi_{a'b'c'}) \label{eq:dphi}\\
|d{*}\varphi|^2(s) &= 5 \cdot g^{00}(s) \cdot g^{aa'} g^{bb'} g^{cc'} g^{dd'} \cdot (\partial_s \psi_{abcd})(\partial_s \psi_{a'b'c'd'}) \label{eq:dstarphi}
\end{align}

The factor $g^{00}$ accounts for the seam-direction metric component. The factors 4 and 5
arise from the antisymmetry of $\varphi$ (3-form) and $\psi$ (4-form) respectively.

\textbf{$C^0$ (supremum) norm:}
\begin{equation}\label{eq:torsion_c0}
\norm{T}_{C^0} := \sup_{s \in [-2,3]} \sqrt{|d\varphi|^2(s) + |d{*}\varphi|^2(s)}
\end{equation}

Evaluated at $N = 100$ Chebyshev nodes on $[0,1]$. The certified bound uses the Chebyshev
coefficient sum $\norm{T^2}_\infty \leq \sum|a_k|$, which is grid-free (tightness $1.12\times$).

\textbf{$L^2$ (integrated) norm:}
\begin{equation}\label{eq:torsion_l2}
\norm{T}_{L^2} := \sqrt{\int_{-2}^{3} \bigl(|d\varphi|^2 + |d{*}\varphi|^2\bigr) \sqrt{\det(g)}\, ds}
\end{equation}

Computed via Clenshaw--Curtis quadrature at Chebyshev nodes.

\subsection{Newton--Kantorovich Norm}\label{sec:nk_norm}

The NK iteration operates on the Banach space $C^0([-2, 3], \Sym^+_7(\mathbb{R}))$ of continuous
maps from the extended neck to the cone of $7 \times 7$ symmetric positive-definite matrices,
equipped with the relative Frobenius supremum norm (\S\ref{sec:metric_dist}). The torsion map
$T : g \mapsto (d\varphi[g], d{*}\varphi[g])$ is viewed as an operator on this space.

The NK contraction parameter is:
\begin{equation}\label{eq:nk_h}
h = \beta \cdot \eta \cdot \omega, \qquad \beta = 1/\lambda_1^\perp, \qquad \eta = \norm{T(g_5)}_{C^0}
\end{equation}
where:
\begin{itemize}
    \item $\beta = 1/\lambda_1^\perp$, with $\lambda_1^\perp = 33.771$ the first nonzero eigenvalue of the linearized
          torsion operator $DF(g_5)$, computed by spectral discretization on the Chebyshev grid
          ($N = 100$ nodes) with Dirichlet-to-Neumann boundary conditions from the ACyl tails.
          The eigenvalue is stable under grid refinement ($N = 200$: same value to 4 digits).
    \item $\eta$ is the certified residual torsion, bounded grid-free via Chebyshev coefficient
          sums: $\norm{T^2}_\infty \leq \sum|a_k|$ (using $|T_k(x)| \leq 1$). Tightness ratio: $1.12\times$.
    \item $\omega$ is the Lipschitz constant of $DF$, computed from the Fr\'echet derivative of the
          torsion operator via centered finite-difference Jacobians at 10 perturbed metrics,
          with a $3\times$ safety factor on the maximum observed value. The certified value is
          $\omega = 0.0713$.
\end{itemize}

\subsection{K3 Verification Norms}\label{sec:k3_norms}

For the $y$-dependent torsion field $T(s, y)$, $y \in \mathrm{K3}$:
\begin{equation}\label{eq:k3_norm}
\norm{T}_{C^0}^{(s,y)} := \sup_{s,y} \sqrt{|d\varphi|^2(s,y) + |d{*}\varphi|^2(s,y)}
\end{equation}
evaluated at $N_s = 100$ seam nodes $\times$ all $N_{\mathrm{K3}} = 220{,}000$ K3 sample points from the
cymyc neural network approximation of the K3 Ricci-flat metric ($\sigma_{\mathrm{metric}} = 0.011$;
see Appendix~\ref{app:k3}).

This norm is used in \S\ref{sec:k3_fiber} to certify that the K3 fiber variation does not increase
the torsion beyond the seam-ansatz bounds.

% ============================================
\section{Torsion Analysis}\label{sec:torsion}
% ============================================

\subsection{Initial Approximation}\label{sec:initial_torsion}

The baseline metric (from Chebyshev fit of a numerically optimized solution) has:

\begin{table}[H]
\centering
\begin{tabular}{ll}
\toprule
\textbf{Quantity} & \textbf{Value} \\
\midrule
$\norm{T}_{C^0}$ & $8.936 \times 10^{-2}$ \\
$\norm{d\varphi}_{C^0}$ & $3.652 \times 10^{-2}$ \\
$\norm{d{*}\varphi}_{C^0}$ & $8.158 \times 10^{-2}$ \\
$|d\varphi|^2 / |d{*}\varphi|^2$ & $0.200000$ (exact $1/5$) \\
Torsion class & $W_2 \oplus W_3$ (99.6\% in $\tau_3$, $\tau_1 \approx 0$) \\
\bottomrule
\end{tabular}
\end{table}

The ratio $|d\varphi|^2/|d{*}\varphi|^2 = 1/5$ is a consequence of $\Gtwo$ representation theory: in the
Fern\'andez--Gray classification~\cite{fernandez1982}, the torsion of the 1D ansatz lies in the
$W_3 = \mathbb{R}^7$ component, where the branching rule $\Lambda^4_7 \leftrightarrow \Lambda^5_7$ gives the $1{:}5$ ratio
between the closure and co-closure contributions.

\subsection{K3 Fiber Verification}\label{sec:k3_fiber}

The 1D seam ansatz treats the K3 fiber metric as spatially homogeneous ($y$-independent
at each $s$). This does \textit{not} assume K3 is flat: the K3 surface has $\chi = 24$ and
nontrivial Riemann curvature, with an intrinsic metric variation of ${\sim}70\%$ across the
surface (Appendix~\ref{app:k3}). The relevant property for the TCS construction is that K3 is
approximately Ricci-flat ($\sigma = 0.011$), which controls the transverse torsion
independently of the seam analysis. What \S\ref{sec:k3_fiber} bounds is the impact of K3 fiber
variation on the \textit{1D seam torsion}: specifically, how much the torsion functional
changes when the constant fiber metric is replaced by the actual $y$-dependent K3 metric,
scaled by the NK correction amplitude.

To certify that the K3 fiber variation does not invalidate the 1D torsion bounds,
we bound the perturbation $T(g + \delta g_{\mathrm{K3}})$ through three quantities:

\begin{enumerate}
    \item \textbf{Fr\'echet derivative.} The operator norm $\norm{\partial T / \partial g_{\mathrm{K3}}}$ of the torsion with
          respect to the 10 independent K3 metric entries (indices $\{2, 3, 4, 5\}$ in the
          symmetric $7 \times 7$ matrix) is evaluated via centered finite differences and bounded
          by Chebyshev coefficient sums (grid-free). Result: $\norm{DT} = 1.21 \times 10^{-5}$.

    \item \textbf{K3 fiber supremum.} Over all 220,000 K3 sample points from the cymyc metric
          (Appendix~\ref{app:k3}), the maximal fiber perturbation is:
          \begin{equation}\label{eq:k3_sup}
          \sup_y \norm{\delta g_{\mathrm{K3}}(y)}_F = 1.80 \times 10^{-3}
          \end{equation}
          where $\delta g_{\mathrm{K3}}(y) = \varepsilon \cdot \delta\hat{R}(y)$ with NK amplitude $\varepsilon = 3.69 \times 10^{-4}$ (from the NK
          convergence ball, \S\ref{sec:nk_cert}) and $\delta\hat{R}(y)$ the K3 shape modes normalized to unit RMS.
          The smallness of $\norm{\delta g_{\mathrm{K3}}}$ despite the large intrinsic K3 variation (${\sim}70\%$) is
          due to the NK amplitude $\varepsilon$, which scales the K3 shape by the size of the metric
          correction needed for torsion-free convergence.

    \item \textbf{Perturbation bound.} By linearity of the leading-order correction,
          \begin{equation}\label{eq:k3_pert}
          \norm{T(g + \delta g_{\mathrm{K3}})} \leq \norm{T(g)} + \norm{DT} \cdot \norm{\delta g_{\mathrm{K3}}} + O(\norm{\delta g_{\mathrm{K3}}}^2)
          \end{equation}
\end{enumerate}

\begin{table}[H]
\centering
\begin{tabular}{ll}
\toprule
\textbf{Quantity} & \textbf{Value} \\
\midrule
Linear K3 correction to torsion & $2.17 \times 10^{-8}$ \\
Quadratic K3 correction & $4.95 \times 10^{-12}$ \\
K3-inclusive certified torsion $\norm{T}_{C^0}^{\mathrm{K3}}$ & $3.154 \times 10^{-5}$ \\
Fraction of 1D torsion & 0.07\% \\
\bottomrule
\end{tabular}
\end{table}

\textbf{Cross-validation}: at the worst K3 point (index 151,325 of 220,000), the actual
torsion ($2.983 \times 10^{-5}$) is \textit{lower} than the 1D value: the linear bound is
conservative. The ratio $|d\varphi|^2/|d{*}\varphi|^2 = 0.200000$ is preserved exactly at every K3
sample point.

\subsection{One-Dimensional Torsion Floor}\label{sec:torsion_floor}

With fixed boundary conditions $g(0) \neq g(1)$ (matching the two ACyl ends), the 1D seam
ansatz has a hard torsion floor:
\begin{equation}\label{eq:torsion_floor}
\min_{g(s),\, g(0)=g_0,\, g(1)=g_1} \norm{T[g]}_{C^0} \approx 0.079
\end{equation}

This floor is \textbf{DOF-independent}: increasing the Chebyshev order $K$ from 5 to 20, or
using direct node-value optimization with 2744 parameters, yields the same floor to
within 1\%. Unconstrained optimization (boundaries free) reaches $\norm{T} = 9.4 \times 10^{-5}$,
confirming the floor originates from the boundary mismatch, not spectral limitations.

\textbf{Transverse correction analysis:} $T^2$ Fourier modes cannot break this floor
(gradient $\partial\mathcal{L}/\partial\alpha = 0$ at $\alpha = 0$ to machine precision). Linearized decomposition shows
100\% of active torsion requires K3 directions (98.1\% of the NK correction is in the
K3 block of the metric).

\subsection{Gauss--Newton Torsion Reduction}\label{sec:gauss_newton}

Removing the fixed boundary constraint and replacing it with Newton--Kantorovich ball
projection ($\norm{\delta g}/\norm{g} \leq 3.76 \times 10^{-4}$ per step), the LBFGS optimizer drives the metric
toward the torsion-free solution. The iteration converges in 5 steps:

\begin{table}[H]
\centering
\begin{tabular}{cccc}
\toprule
\textbf{Iteration} & $\norm{T}_{C^0}$ & \textbf{Cumulative reduction} & \textbf{Cumulative} $\delta g / g$ \\
\midrule
0 & $8.936 \times 10^{-2}$ & $1.0\times$ & 0 \\
1 & $6.364 \times 10^{-2}$ & $1.40\times$ & 0.036\% \\
2 & $3.791 \times 10^{-2}$ & $2.36\times$ & 0.072\% \\
3 & $1.220 \times 10^{-2}$ & $7.33\times$ & 0.107\% \\
4 & $2.626 \times 10^{-5}$ & $3403\times$ & 0.124\% \\
5 & $2.485 \times 10^{-5}$ & $\mathbf{3596\times}$ & 0.124\% \\
\bottomrule
\end{tabular}
\end{table}

Steps 1--3 exhibit linear (geometric) convergence; step~4 enters the superlinear
(Newton) regime with a factor ${\sim}470$ reduction in a single step.

\begin{figure}[ht]
\centering
\includegraphics[width=0.7\textwidth]{figures/fig2_torsion_convergence.pdf}
\caption{Convergence of the torsion norm $\norm{T}_{C^0}$ over 5 Gauss--Newton
iterations (log scale). Steps 0--3 show linear convergence; step~4 enters the
superlinear (Newton) regime with a single-step factor of ${\sim}470\times$. The final
NK-certified distance to the torsion-free metric $g^*$ is annotated.}
\label{fig:torsion_convergence}
\end{figure}

\textit{Note.} The iteration table above corresponds to the LBFGS path reproduced by the
companion notebook (\S\ref{sec:notebook}), which is the verifiable artifact. The embedded optimized
coefficients (from the original computation) have $\norm{T}_{C^0} = 2.984 \times 10^{-5}$; the
companion's independent re-run reaches $2.485 \times 10^{-5}$. Both are consistent with the
NK certification, which only requires $\norm{T} < \varepsilon_0 = 0.1$.

\textbf{Separate closure and co-closure at the final iterate:}

\begin{table}[H]
\centering
\begin{tabular}{ccccc}
\toprule
\textbf{Norm} & $\norm{d\varphi}$ & $\norm{d{*}\varphi}$ & $\norm{T}$ & $|d\varphi|^2/|d{*}\varphi|^2$ \\
\midrule
$C^0$ & $1.218 \times 10^{-5}$ & $2.724 \times 10^{-5}$ & $2.984 \times 10^{-5}$ & $0.200000$ \\
\bottomrule
\end{tabular}
\end{table}

These values correspond to the embedded optimized coefficients (\texttt{OPT\_COEFFS} in the
companion notebook). The companion's independent re-run achieves
$\norm{T}_{C^0} = 2.485 \times 10^{-5}$ with the same structural properties.

Both $d\varphi$ and $d{*}\varphi$ are individually small: the near-vanishing of $\norm{T}$ is not due to
cancellation between the two terms.

% ============================================
\section{Certification}\label{sec:certification}
% ============================================

\subsection{Newton--Kantorovich Convergence Certificate}\label{sec:nk_cert}

\textbf{Setting.} Let $X = C^0([-2, 3], \Sym^+_7(\mathbb{R}))$ with the relative Frobenius supremum
norm (\S\ref{sec:metric_dist}), and let $F : X \to C^0([-2, 3], \Lambda^4 \oplus \Lambda^5)$ be the torsion map
$F(g) = (d\varphi[g], d{*}\varphi[g])$.

\begin{theorem}[NK certification]\label{thm:nk}
Let $g_5 \in X$ be the optimized 169-parameter
Chebyshev--Cholesky metric (\S\ref{sec:gauss_newton}). Suppose:
\begin{enumerate}
    \item The linearization $DF(g_5)$ has a bounded inverse with $\norm{DF(g_5)^{-1}} \leq \beta$.
    \item The residual satisfies $\eta := \norm{F(g_5)}_{C^0} < \infty$.
    \item $DF$ is Lipschitz continuous in $B(g_5, r)$ with constant $\omega$.
    \item The product $h := \beta \cdot \eta \cdot \omega < 1/2$.
\end{enumerate}
Then there exists a unique $g^* \in B(g_5, r^*)$ with $F(g^*) = 0$.
\end{theorem}

\textbf{Verification of hypotheses.} All constants are computed with conservative safety
factors (spectral discretization for $\beta$, finite-difference Jacobians for $\omega$):

\begin{table}[H]
\centering
\begin{tabular}{p{2.8cm}p{8cm}l}
\toprule
\textbf{Quantity} & \textbf{How computed} & \textbf{Value} \\
\midrule
$\beta = 1/\lambda_1^\perp$ & Spectral discretization of $DF(g_5)$, $N = 100$ Chebyshev nodes; stable to $N = 200$ & $0.02961$ \\[0.3em]
$\eta = \norm{F(g_5)}_{C^0}$ & Grid-free Chebyshev coefficient bound $\sum|a_k|$ (tightness $1.12\times$) & $3.152 \times 10^{-5}$ \\[0.3em]
$\omega = \Lip(DF)$ & Fr\'echet derivative via FD Jacobian at 10 perturbed metrics, $3\times$ safety & $0.0713$ \\[0.3em]
$h = \beta \cdot \eta \cdot \omega$ & Product & $\mathbf{6.65 \times 10^{-8}}$ \\
\bottomrule
\end{tabular}
\end{table}

Since $h = 6.65 \times 10^{-8} < 1/2$, with a margin of $\mathbf{\times 7.5}$~\textbf{million}, the hypotheses
are satisfied. The maximum $\omega$ that would still permit certification is $535{,}732$: a
factor $7{,}500$ above the computed value, demonstrating robustness against any
reasonable tightening of the Lipschitz bound.

\textbf{K3-inclusive bound.} The torsion bound $\eta$ includes the K3 fiber contribution
via the perturbation analysis of \S\ref{sec:k3_fiber}: the K3-inclusive torsion bound
$\norm{T}_{C^0}^{\mathrm{K3}} = 3.154 \times 10^{-5}$ (over all 220,000 K3 sample points) differs
from the 1D bound by 0.07\%. The NK parameter $h$ remains far below $1/2$.

\textbf{Certified distance to} $g^*$\textbf{:}
\begin{equation}\label{eq:dist_gstar}
\dist(g_5, g^*) \leq 4.86 \times 10^{-6} \qquad \text{(relative Frobenius norm)}
\end{equation}

Within the computational precision of this certificate, the torsion-free metric $g^*$
exists and is unique in a ball around the explicit iterate $g_5$. We emphasize that this
is a \textit{computational} NK certificate: the constants $\beta$, $\omega$ are obtained from spectral
discretization and finite-difference Jacobians (with a $3\times$ safety factor), not from
rigorous interval arithmetic. The margin of ${\times}7.5$~million on the NK parameter $h$
provides substantial robustness against numerical error, but full interval
verification of operator norms is left for future work.

\subsection{Interval Arithmetic Verification}\label{sec:interval}

Grid-free certification using interval Chebyshev arithmetic:

\begin{table}[H]
\centering
\begin{tabular}{lll}
\toprule
\textbf{Property} & \textbf{Result} & \textbf{Method} \\
\midrule
$\det(g) = 65/32$ & \textbf{Exact} (algebraic, by construction) & $\alpha^7 \det(L) = \sqrt{65/32}$ \\
$g = LL^\top$, $L_{ii} > 0$ & \textbf{Structural} (softplus guarantees) & $L_{ii} = \log(1+\exp(\cdot)) > 0$ \\
$\lambda_{\min}(g) \geq 0.373$ & Certified via Gershgorin & Chebyshev coefficient intervals \\
$\kappa(g) \leq 3.88$ & Certified condition number & Same \\
Metric variation $< 0.16\%$ & Along neck & Chebyshev bound propagation \\
\bottomrule
\end{tabular}
\end{table}

\subsection{Holonomy Proof}\label{sec:holonomy}

\begin{theorem}\label{thm:holonomy}
The holonomy group of the certified metric $g^*$ on $\Kseven$ is exactly $\Gtwo$.
\end{theorem}

\begin{proof}
By Joyce~\cite{joyce2000}, Proposition~11.2.3: on a compact 7-manifold $M$ with a
torsion-free $\Gtwo$-structure and $\pi_1(M) = \{1\}$, the holonomy is $\Hol(g) = \Gtwo$ if and only
if $b_1(M) = 0$.

Verification of the hypotheses:
\begin{enumerate}
    \item \textbf{Compactness}: $\Kseven$ is compact by the TCS construction (gluing of compact pieces
          with finite overlap).
    \item \textbf{Torsion-free}: The metric $g^*$ is torsion-free by the NK certificate (\S\ref{sec:nk_cert}).
    \item \textbf{Simply connected}: $\pi_1(\Kseven) = \{1\}$ by the argument in \S\ref{sec:topology}.
    \item $b_1 = 0$: Follows from simple connectivity.
\end{enumerate}
Therefore $\Hol(g^*) = \Gtwo$.
\end{proof}

\textbf{Numerical indicators} (at the optimized iterate $g_5$, supporting the formal proof):
\begin{itemize}
    \item $|\varphi|^2 = 42.0000000000$ (error: $2.1 \times 10^{-14}$), $\Gtwo$ calibration preserved
    \item $\norm{\mathrm{Rm}}_{C^0} = 4.20 \times 10^{-5}$, metric is non-flat (curvature nonzero)
    \item $\norm{\Gamma}_{C^0} = 2.81 \times 10^{-6}$, neck nearly flat (curvature concentrated in ACyl bulks)
\end{itemize}

% ============================================
\section{Summary of Geometric Invariants}\label{sec:invariants}
% ============================================

\begin{table}[H]
\centering
\begin{tabular}{lll}
\toprule
\textbf{Invariant} & \textbf{Value} & \textbf{Status} \\
\midrule
$\det(g)$ & $65/32$ (exact) & Algebraic by construction \\
$|\varphi|^2$ & 42 (error $< 10^{-14}$) & $\Gtwo$ calibration \\
$|\psi|^2$ & 168 & Coassociative calibration \\
$|d\varphi|^2/|d{*}\varphi|^2$ & $1/5$ (exact to $10^{-10}$) & $\Gtwo$ representation theory~\cite{fernandez1982} \\
$\tau_1$ & $< 10^{-9}$ & Torsion class = $W_2 \oplus W_3$ \\
$\tau_3$ fraction & 99.6\% & Initial approximation \\
$\norm{T}_{C^0}$ (final) & $2.984 \times 10^{-5}$ & After 5 iterations \\
$\dist(g_5, g^*)$ & $\leq 4.86 \times 10^{-6}$ & Computational NK certificate (\S\ref{sec:nk_cert}) \\
$\lambda_{\min}(g)$ & 0.822 & SPD certified \\
$\kappa(g)$ & $\leq 3.88$ & Well-conditioned \\
$b_2$ & 21 & \S\ref{sec:topology} \\
$b_3$ & 77 & \S\ref{sec:topology} \\
$\Hol(g^*)$ & $\Gtwo$ & \S\ref{sec:holonomy} \\
\bottomrule
\end{tabular}
\end{table}

% ============================================
\section{Discussion}\label{sec:discussion}
% ============================================

\subsection{Limitations}\label{sec:limitations}

The principal limitation of this work is the \textbf{seam ansatz}: the metric $g(s)$ depends
only on the neck coordinate $s$, with constant fiber metrics at each cross-section. The
full metric on $\Kseven$ would be a field $g(s, \theta, y, \psi)$ depending on all 7 coordinates.
This paper does \textit{not} solve the full 7D torsion-free equation $d\varphi = 0$, $d{*}\varphi = 0$ on
the compact manifold. The K3 fiber verification (\S\ref{sec:k3_fiber}) bounds the fiber variation
\textit{a posteriori} ($\leq 0.07\%$ of torsion), and the NK certificate (\S\ref{sec:nk_cert}) establishes
convergence within the 1D family. Extension to the full fiber-coupled problem
remains the main open direction (\S\ref{sec:future}).

A second limitation is that the ACyl extension (\S\ref{sec:acyl}) uses a simple exponential decay model. The true
asymptotic behavior involves the spectrum of the Laplacian on the K3 cross-section
and higher-order corrections. The $97.2\% / 2.8\%$ torsion split between neck and tails
suggests this is a minor effect, but it remains uncontrolled.

\subsection{Comparison with Existing Work}\label{sec:comparison}

Previous numerical work on $\Gtwo$ metrics has focused on non-compact examples, particularly
cohomogeneity-one metrics on bundles over $S^3$ and $\mathbb{CP}^2$~\cite{brandhuber2001}, where the ODE structure
permits high-precision computation. For compact $\Gtwo$ manifolds, the existence results
of Joyce~\cite{joyce1996, joyce2000} and the TCS construction of Kovalev~\cite{kovalev2003} and CHNP~\cite{chnp2015} establish metrics
perturbatively but do not yield explicit numerical values.

Concurrent work by Heyes, Hirst, S\'a Earp and Silva~\cite{heyes2026} applies neural networks to
approximate $\Gtwo$-structures on contact Calabi--Yau 7-manifolds (links in $S^9$), achieving
learned 3-forms and torsion estimates via data-driven methods. Their approach and ours
are complementary: \cite{heyes2026} works on non-compact (cone/link) geometries with a neural
architecture trained on point samples, while the present paper works on a compact TCS
manifold with an analytical (Chebyshev) parametrization and a convergence certificate.

\subsection{Future Directions}\label{sec:future}

\begin{enumerate}
    \item \textbf{Full fiber-coupled metric.} Extend the Chebyshev--Cholesky parametrization to
          include $y$-dependence on the K3 fiber, using the cymyc K3 metric (Appendix~\ref{app:k3}) as
          initial data. The K3 correction analysis (\S\ref{sec:k3_fiber}, \S\ref{sec:torsion_floor}) provides a starting point.

    \item \textbf{Spectral geometry.} With an explicit metric available, Laplacian eigenvalues and
          harmonic forms can be computed numerically, quantities not yet available for
          compact $\Gtwo$ manifolds.

    \item \textbf{Other topological types.} Apply the same pipeline to other TCS manifolds from the
          CHNP classification~\cite{chnp2015}, to understand how the metric structure depends on $(b_2, b_3)$.

    \item \textbf{Rigorous interval verification.} Upgrade the computational NK certificate to a
          computer-assisted proof by replacing the finite-difference Lipschitz bound $\omega$ and
          the spectral eigenvalue $\beta$ with rigorous interval arithmetic bounds, in the spirit
          of validated numerics for PDEs.

    \item \textbf{Comparison with flow methods.} Compare the Chebyshev--Cholesky metric with results
          from Laplacian flow~\cite{lotay2019} or Hitchin flow, which provide alternative approaches to $\Gtwo$
          metrics.
\end{enumerate}

% ============================================
\section{Reproducibility}\label{sec:reproducibility}
% ============================================

\subsection{Supplementary Data}\label{sec:supp_data}

The complete metric is distributed as a single JSON file containing:

\begin{table}[H]
\centering
\begin{tabular}{ll}
\toprule
\textbf{Field} & \textbf{Content} \\
\midrule
\texttt{chebyshev\_coefficients} & $6 \times 28$ coefficient matrix $C$ (float64) \\
\texttt{acyl\_decay.gamma} & Decay rate $\gamma = 5.811297$ \\
\texttt{G0\_star} & Baseline $7 \times 7$ metric at neck midpoint \\
\texttt{metadata} & Betti numbers, determinant, parameter counts \\
\bottomrule
\end{tabular}
\end{table}

This file, together with the reconstruction algorithm (\S\ref{sec:reconstruction}), fully specifies the
metric at any point of the extended neck domain.

\subsection{Companion Notebook}\label{sec:notebook}

A self-contained Jupyter notebook is provided as supplementary material. It requires
only PyTorch and NumPy (no external dependencies), embeds all 169 parameters, and
executes the full verification chain:

\begin{enumerate}
    \item Metric reconstruction from Chebyshev coefficients
    \item Structural checks (SPD, determinant, calibration norms)
    \item Torsion computation and Joyce iteration (5 steps)
    \item Newton--Kantorovich certification
    \item SHA-256 hashes of all embedded coefficient tensors
    \item JSON export of all results with UTC timestamp
\end{enumerate}

Total runtime: under a minute on CPU. The notebook outputs a timestamped JSON
manifest (\texttt{G2\_Metric\_Companion\_results.json}) containing all check results,
coefficient hashes, and environment metadata for tamper-evident reproducibility.

\textbf{Expected output:} 15/15 checks pass. Key values: $\norm{T}_{C^0} = 2.98 \times 10^{-5}$,
$h = 6.65 \times 10^{-8}$, $|\varphi|^2 = 42$ (to 14 digits), $\det(g) = 65/32$ (to 15 digits).

\subsection{Computational Environment}\label{sec:environment}

\begin{itemize}
    \item Python 3.x with PyTorch (float64 throughout, IEEE 754 double precision)
    \item Runs on CPU or any CUDA device (tested on NVIDIA RTX 2050 and Google Colab CPU)
    \item K3 fiber metric computed using cymyc~\cite{cymyc} on a cloud GPU (72~min)
    \item No dependencies beyond PyTorch and NumPy for verification
\end{itemize}

% ============================================
% AUTHOR'S NOTE
% ============================================

\vspace{1em}
\noindent\rule{\textwidth}{0.2pt}

\subsection*{Author's note}

This work was developed through sustained collaboration between the author and several AI systems, primarily Claude (Anthropic), with contributions from GPT (OpenAI) for specific mathematical insights. The architectural decisions and many key derivations emerged from iterative dialogue sessions over several months. This collaboration follows a transparent crediting approach for AI-assisted mathematical research. The value of any proposal depends on mathematical coherence and empirical accuracy, not origin. Mathematics is evaluated on results, not r\'esum\'es.

\vspace{1em}
\noindent\rule{\textwidth}{0.2pt}

% ============================================
% REFERENCES
% ============================================

\begin{thebibliography}{99}

\bibitem{harvey1982}
R.~Harvey and H.\,B.~Lawson,
``Calibrated geometries,''
\textit{Acta Math.} \textbf{148} (1982), 47--157.

\bibitem{bryant1987}
R.\,L.~Bryant,
``Metrics with exceptional holonomy,''
\textit{Ann.\ Math.} \textbf{126}(3) (1987), 525--576.

\bibitem{joyce1996}
D.\,D.~Joyce,
``Compact Riemannian 7-manifolds with holonomy $\mathrm{G}_2$.\ I, II,''
\textit{J.\ Diff.\ Geom.} \textbf{43}(2) (1996), 291--328 and 329--375.

\bibitem{joyce2000}
D.\,D.~Joyce,
\textit{Compact Manifolds with Special Holonomy},
Oxford University Press, 2000.

\bibitem{kovalev2003}
A.\,G.~Kovalev,
``Twisted connected sums and special Riemannian holonomy,''
\textit{J.\ Reine Angew.\ Math.} \textbf{565} (2003), 125--160.

\bibitem{chnp2015}
A.~Corti, M.~Haskins, J.~Nordstr\"om and T.~Pacini,
``$\mathrm{G}_2$-manifolds and associative submanifolds via semi-Fano 3-folds,''
\textit{Duke Math.\ J.} \textbf{164}(10) (2015), 1971--2092.

\bibitem{fernandez1982}
M.~Fern\'andez and A.~Gray,
``Riemannian manifolds with structure group $\mathrm{G}_2$,''
\textit{Ann.\ Mat.\ Pura Appl.} \textbf{132} (1982), 19--45.

\bibitem{cymyc}
M.~Larfors, A.~Lukas and F.~Ruehle,
``Calabi--Yau metrics from machine learning,''
\textit{JHEP} \textbf{2022}, 232.
(\texttt{cymyc} software package.)

\bibitem{lotay2019}
J.\,D.~Lotay and Y.~Wei,
``Laplacian flow for closed $\mathrm{G}_2$ structures:
Shi-type estimates, uniqueness and compactness,''
\textit{Geom.\ Funct.\ Anal.} \textbf{29} (2019), 1048--1110.

\bibitem{brandhuber2001}
A.~Brandhuber, J.~Gomis, S.\,S.~Gubser and S.~Gukov,
``Gauge theory at large $N$ and new $\mathrm{G}_2$ holonomy metrics,''
\textit{Nuclear Phys.\ B} \textbf{611} (2001), 179--204.

\bibitem{heyes2026}
E.~Heyes, E.~Hirst, H.\,N.~S\'a Earp and T.\,S.\,R.~Silva,
``Neural and numerical methods for $\mathrm{G}_2$-structures on contact Calabi--Yau 7-manifolds,''
arXiv:2602.12438, 2026.

\end{thebibliography}

% ============================================
% APPENDICES
% ============================================

\appendix

\section{The Baseline Metric ($7 \times 7$)}\label{app:baseline}

The baseline metric $g_0$ at the neck midpoint ($s = 0.5$) is approximately:
\[
g_0 \approx \mathrm{diag}(6.789,\; 2.904,\; 1.093,\; 1.094,\; 1.092,\; 1.095,\; 2.904)
\]
with off-diagonal elements $< 0.007$ (the metric is nearly diagonal). The precise
values are stored in the supplementary data file under the key \texttt{G0\_star}.

The seam direction (index~0) has the largest metric component (${\sim}6.8$), reflecting the
extended neck geometry. The K3 directions (indices 2--5) are approximately unit scale
(${\sim}1.09$), and the $T^2$ directions (indices 1, 6) are intermediate (${\sim}2.9$). This eigenvalue
structure is consistent with the index assignment: K3 = $\{2, 3, 4, 5\}$ (4 directions,
uniform scale), $T^2 = \{1, 6\}$ (2 circle directions, equal scale).

\section{K3 Fiber Metric}\label{app:k3}

The K3 fiber metric is computed using the cymyc neural network~\cite{cymyc} trained on the
complete intersection $\mathrm{CI}(1,2,2,2) \subset \mathbb{P}^6$:

\begin{table}[H]
\centering
\begin{tabular}{ll}
\toprule
\textbf{Parameter} & \textbf{Value} \\
\midrule
Sample points & 220,000 \\
Metric accuracy $\sigma$ & 0.011 \\
SPD fraction & 100\% (220,000/220,000) \\
Hyperk\"ahler triple & $J_I \cdot J_J = J_K$ to $2.2 \times 10^{-10}$ \\
HK norms & $\Tr(\omega\, g^{-1}\, \omega\, g^{-1}) = -4.00$ for all three \\
Intrinsic K3 variation & 69.8\% ($\sigma/\norm{g}$) \\
NK-scaled variation in $\Gtwo$ & 0.037\% \\
\bottomrule
\end{tabular}
\end{table}

\textbf{Clarification on the two variation measures.} The ``intrinsic K3 variation'' (69.8\%)
measures how much the Ricci-flat K3 metric $g_{\mathrm{CY}}(y)$ varies across the K3 surface:
$\sigma(g_{\mathrm{CY}})/\langle\norm{g_{\mathrm{CY}}}\rangle$. This is large and \textit{expected}: K3 has Euler characteristic $\chi = 24$
and nontrivial Riemann curvature; a constant (flat) metric on K3 does not exist.
The seam ansatz does not assume K3 is flat; it uses a spatially homogeneous
($y$-independent) fiber metric at each value of $s$, representing the average fiber
geometry.

The ``NK-scaled variation'' (0.037\%) is a different quantity: the NK amplitude
$\varepsilon = 3.69 \times 10^{-4}$ (the size of the metric correction in the NK iteration) multiplied
by the intrinsic variation. This is the quantity that enters the torsion perturbation
bound (\S\ref{sec:k3_fiber}). The two measures are related by:
\[
\text{NK-scaled} = \varepsilon \times \text{intrinsic} = 3.69 \times 10^{-4} \times 69.8\% \approx 0.037\%
\]

The relevant measure for the TCS construction is not whether K3 is flat (it is not),
but whether it is Ricci-flat. The cymyc metric accuracy $\sigma = 0.011$ (1.1\%) controls
the transverse torsion contribution; the 69.8\% measures the Riemann curvature
(nonzero, as required by topology).

\section{Chebyshev Coefficients}\label{app:cheb}

The full $6 \times 28$ coefficient matrix is provided in machine-readable form in the
supplementary data file under the key \texttt{chebyshev\_coefficients}.

For reference, the dominant row ($k = 0$, constant mode):

\begin{lstlisting}
c_0 = [2.5592, 0.0001, 1.3751, 0.0000, 0.0003, 1.0452, 0.0000, -0.0002,
       0.0061, 1.0456, 0.0001, 0.0000, -0.0038, 0.0030, 1.0456, 0.0001,
       0.0000, 0.0003, -0.0026, -0.0003, 1.0460, 0.0001, 0.0000, 0.0002,
       -0.0003, 0.0002, -0.0002, 1.3751]
\end{lstlisting}

(4 significant figures; full 16-digit precision in the supplementary file.)

The 7 diagonal entries of $L$ at $k = 0$ are: $\{2.559, 1.375, 1.045, 1.046, 1.046, 1.046, 1.375\}$,
which after softplus give: $\{2.637, 1.639, 1.366, 1.367, 1.367, 1.367, 1.639\}$. These determine
the scales of the 7 fiber directions: seam $\approx \sqrt{6.8}$, $T^2 \approx \sqrt{2.9}$, K3 $\approx \sqrt{1.09}$.

\section{Remark on Related Work}\label{app:related}

The specific topological type ($b_2 = 21$, $b_3 = 77$) and the determinant normalization
$\det(g) = 65/32$ used in this paper are motivated by the Geometric Information Field
Theory (GIFT) framework (de La Fourni\`ere, 2026, technical report), which studies
M-theory compactifications on $\Gtwo$ manifolds. In that context, the building block choice
and Betti numbers are selected to satisfy anomaly cancellation constraints related to
$\E_8 \times \E_8$ gauge structure, and the determinant value arises from a topological formula
involving the dimensions of $\Gtwo$ and $\E_8$.

The present paper makes no physical claims and treats these as input data for a purely
geometric construction. In particular, the mathematical results (torsion bounds,
Newton--Kantorovich certification, and the holonomy proof) hold for any compact TCS
7-manifold admitting a $\Gtwo$-structure with sufficiently small torsion, regardless of
physical interpretation.

Concurrent work by Heyes, Hirst, S\'a Earp and Silva~\cite{heyes2026} applies neural networks to
$\Gtwo$-structures on contact Calabi--Yau 7-manifolds, with cross-citation of our
preliminary numerical results. The two approaches are complementary (see \S\ref{sec:comparison}).

\vspace{2em}
\begin{center}
\textit{Manuscript prepared March 2026.}
\end{center}

\end{document}
